%% SYNTHETIC TEST UNIT generated by writing/tools/make_outline_book.py from FINAL_OUTLINE_v2 -- not book text
\chapter{Wilson 线：Feynman 规则、程函恒等式与尖点}\label{app:appB}
本单元为测试用的合成文本（由 FINAL\_OUTLINE\_v2 自动生成），陈述内容从略。

\section{第 1 节（测试）}\label{sec:appB-definition}

\begin{proposition}[Wilson 线的 Feynman 规则 \tagR]\label{prop:appB-rules}
（$R_{\mathrm{def}}$：由算符定义直接推出。）
（测试用合成陈述：内容从略。）
\uses{def:prelim-wilson}
\end{proposition}

\section{第 2 节（测试）}\label{sec:appB-rules}

\begin{lemma}[程函恒等式（洗牌积）]\label{lem:appB-shuffle}
（测试用合成陈述：内容从略。）
% uses: none
\end{lemma}

\section{第 3 节（测试）}\label{sec:appB-eikonal}

\begin{proposition}[非阿贝尔指数化 \tagA]\label{prop:appB-webs}
（测试用合成陈述：内容从略。）
\uses{prop:appB-rules, lem:appB-shuffle}
\end{proposition}
\begin{proofsketch}测试。\end{proofsketch}
\paragraph{地位与缺口}测试：本陈述的地位由\cref{prop:sudakov-webs}决定。

\section{第 4 节（测试）}\label{sec:appB-cusp}

\begin{proposition}[单圈尖点反常维度 \tagR]\label{prop:appB-cusp-oneloop}
（$R_{\mathrm{PT}}$：在重整化微扰论中对红外正规化的部分子态严格。）
（测试用合成陈述：内容从略。）
\uses{prop:appB-rules, def:prelim-dr-conventions, def:prelim-msbar}
\end{proposition}
